\documentclass[../../main.tex]{subfiles}
\begin{document}
% SUPPLEMENT: id=v14-ch04-regular-twins-construction; parent=ch04-notes-separation-history; kind=omitted-proof+construction+motivation+hypotheses; origin=independent-derivation-with-attributed-construction
\begin{remark}\phantomsection\label{supplement:ch04-regular-constant-construction}
\textbf{正则空间仍可能只有常值连续函数：完整构造的准备（重构新增）。}这一存在性不能由平凡拓扑的例子代替，因为本书的正则性包含 $T_1$。以下先造一个正则空间中的两个不同点，使所有连续实函数在它们处同值，再用商空间把这种约束传播到所有点。前一构造采用 Kalemba--Plewik\cite[\S2]{supp-kalemba-plewik-2017}方法的一个具体特例；后一拼接原则是 Baggs\cite[\S2]{supp-baggs-1989}所述的 van Douwen 构造。下文独立展开全部验证，特别补出商空间正则性的证明。
\end{remark}
\begin{definition*}[用于构造的分层空间]\phantomsection\label{definition:ch04-layered-twin-space}
\textbf{补充构造（重构新增）。}取 $D=\mathbb Z\times\omega_1$，其中 $\omega_1$ 为第一不可数序数；令 $D_k=\{k\}\times\omega_1$。在 $P=D\times D$ 中记 $d_x=(x,x)$、$\Delta_k=\{d_x:x\in D_k\}$、$\Delta=\bigcup_k\Delta_k$。

令 $P\setminus\Delta$ 的每个点孤立。对 $x\in D_k$，规定 $d_x$ 的基本邻域为
\begin{align*}
N(x;G,F)=\{d_x\}
\cup\bigl(\{x\}\times(D_{k-1}\setminus G)\bigr)
\cup\bigl((D_{k+1}\setminus F)\times\{x\}\bigr),
\end{align*}
其中 $G\subseteq D_{k-1}$ 有限，$F\subseteq D_{k+1}$ 至多可数。再添两个新点 $a,b$，得 $S=P\cup\{a,b\}$。$P$ 的开集仍开，并规定以下集合分别为 $a,b$ 的邻域基：
\begin{align*}
A_m&=\{a\}\cup\bigcup_{k>m}(D_k\times D),\\
B_m&=\{b\}\cup\left(\bigcup_{k\le m}(D_k\times D)\right)\setminus\Delta_m,
\qquad m\in\mathbb Z.
\end{align*}
这里第二式的差集仅删去 $P$ 中的 $\Delta_m$，不删去 $b$。
\end{definition*}
\begin{proposition}[构造中的两个点]\label{proposition:ch04-regular-twins}
\textbf{补充结论（重构新增）。}$S$ 是 $T_1$ 正则空间，$a\ne b$，而每个连续函数 $f:S\to\mathbb R$ 都满足 $f(a)=f(b)$。
\end{proposition}
\begin{proof}
\textbf{第一步：拓扑与正则性。}固定 $x$ 的两个基本邻域相交，等于将两个有限例外集合、两个可数例外集合分别取并后所得的基本邻域。不同对角点的基本邻域不会含对方的对角点，它们相交处只有孤立点。因此这些集合确实给出 $P$ 的拓扑基。

每个 $N(x;G,F)$ 在 $P$ 中既开又闭。要验证闭性，只需给其他对角点 $d_y$ 一个避开它的基本邻域。若 $y\in D_{k-1}$，可能的交点只有 $(x,y)$，把 $x$ 加入 $d_y$ 的可数例外集合即可排除；若 $y\in D_{k+1}$，可能的交点只有 $(y,x)$，把 $x$ 加入 $d_y$ 的有限例外集合即可排除。其他层以及同层不同的 $y$ 没有可能交点。对外部孤立点，单点邻域已经避开该集。

$A_m\cap P$ 在 $P$ 中开：所含的对角点层数 $k>m$，其基本邻域仅使用第一坐标层 $k,k+1$，仍在该集合中。$B_m\cap P$ 中的对角点层数 $k<m$，其基本邻域第一坐标层 $k,k+1\le m$，且所有分支点均非对角点，所以也包含在 $B_m$ 中。因此添加两点的邻域规则相容，确实给出 $S$ 的拓扑；$A_m$ 随 $m$ 增大而缩小，$B_m$ 随 $m$ 减小而缩小。

各单点在 $S$ 中闭。对角点之外的其他对角点可按上一段排除该单点；若需排除一个指定非对角点，只须把它的相应坐标加入有限或可数例外集。$a,b$ 的邻域可把任意指定的有限层排除；$P$ 为开子空间，而两个新增点也有互不包含对方的邻域。这逐点验证了每个单点的补集开，即 $T_1$。

前述 $N(x;G,F)$ 的第一坐标仅涉及两层，故两个新点都有避开它的邻域；它在 $S$ 中仍闭。孤立单点也由 $T_1$ 而闭。最后直接核对
\begin{align*}
\overline{A_m}=A_m\cup\Delta_m,
\qquad \overline{B_m}=B_m\cup\Delta_m.
\end{align*}
第一式中，$\Delta_m$ 的每个基本邻域都有未被可数例外删光的 $D_{m+1}$ 竖直分支，故遇到 $A_m$；更低层对角点有完全避开的基本邻域，外部孤立点不在闭包中，$B_m$ 则把 $b$ 与 $A_m$ 分开。第二式中，$\Delta_m$ 的每个基本邻域都有未被有限例外删光的 $D_{m-1}$ 水平分支，故遇到 $B_m$；更高层与 $a$ 都可用相应基本邻域避开。

所以 $\overline{A_{m+1}}\subseteq A_m$，$\overline{B_{m-1}}\subseteq B_m$。每个点的每个开邻域里，都能找到一个开邻域，其闭包仍包含在原邻域中；对 $P$ 中的点用开闭基，对 $a,b$ 用这两个包含关系。这与 $T_1$ 一起证明正则性。

\textbf{第二步：一个闭包传播性质。}本步及下一步中的 $V$ 均为 $P$ 的子集，其闭包在 $P$ 中取。若开集 $V\subseteq P$ 包含 $\Delta_k$ 的一个不可数子集 $H$，则 $\Delta_{k-1}\setminus\overline V$ 有限。反设其无限，从中取两两不同的 $d_{y_j}$，$j\ge1$，并各取避开 $V$ 的基本邻域 $N(y_j;G_j,F_j)$。这里每个 $F_j\subseteq D_k$ 可数，故 $\bigcup_jF_j$ 可数。$H$ 不可数，能取 $d_x\in H$ 使 $x\notin\bigcup_jF_j$。再取 $N(x;G,F)\subseteq V$。$G$ 有限，而 $y_j$ 两两不同，存在 $j$ 使 $y_j\notin G$。此时
\begin{align*}
(x,y_j)\in N(x;G,F)\cap N(y_j;G_j,F_j),
\end{align*}
因为它既在第一个邻域的水平分支中，也在第二个邻域的竖直分支中，矛盾。传播性质得证。

\textbf{第三步：连续函数在每层几乎恒定。}设 $f:S\to\mathbb R$ 连续，先限制到 $P$。固定层 $k$ 及 $\varepsilon>0$，把 $\mathbb R$ 覆盖为可数个长度为 $\varepsilon$ 的闭区间。不可数的 $\Delta_k$ 不可能是可数个可数集的并，所以某区间 $[t,t+\varepsilon]$ 的原像与 $\Delta_k$ 的交不可数。取
\begin{align*}
V=f^{-1}((t-\varepsilon,t+2\varepsilon))\cap P.
\end{align*}
前一步使 $\Delta_{k-1}$ 除有限点外包含于 $\overline V$，而连续性给
\begin{align*}
\overline V\subseteq f^{-1}([t-\varepsilon,t+2\varepsilon])\cap P.
\end{align*}
因此 $f$ 在前一层除有限点外取值于一个长度为 $3\varepsilon$ 的闭区间。

对任意指定层 $k$，用上一结论的第 $k+1$ 层作为起点，分别取 $\varepsilon=1/n$。得到有限例外集 $E_n\subseteq\Delta_k$ 及长度 $3/n$ 的闭区间 $J_n$，使 $f(\Delta_k\setminus E_n)\subseteq J_n$。$\bigcup_nE_n$ 可数，不能覆盖不可数的 $\Delta_k$。取 $p_k$ 在其外，令 $c_k=f(p_k)$。对任何同样在其外的 $p$，$f(p)$ 与 $c_k$ 同属每个 $J_n$，故 $|f(p)-c_k|\le3/n$ 对全部 $n$ 成立，即 $f(p)=c_k$。于是每层除可数点外，$f$ 恒为某个 $c_k$。

这些 $c_k$ 全部相等。任给 $\delta>0$，开集 $f^{-1}((c_k-\delta,c_k+\delta))\cap P$ 包含 $\Delta_k$ 的余可数子集，故第二步使 $\Delta_{k-1}$ 除有限点外在其闭包中。在这层再避开给出 $c_{k-1}$ 的可数例外集，可找到一点同时满足函数值为 $c_{k-1}$ 且落在 $[c_k-\delta,c_k+\delta]$。故 $|c_{k-1}-c_k|\le\delta$；任意 $\delta$ 给出 $c_{k-1}=c_k$。沿整数层传递，记共同值为 $c$。

最后，由 $a$ 处连续性，每个 $\varepsilon>0$ 都有 $A_m$ 包含于 $f^{-1}((f(a)-\varepsilon,f(a)+\varepsilon))$。其中 $\Delta_{m+1}$ 有取值为 $c$ 的点，故 $|c-f(a)|<\varepsilon$。于是 $f(a)=c$。对 $b$ 用 $B_m$ 和 $\Delta_{m-1}$ 同理得 $f(b)=c$，证明两个点无法被连续实函数区分。
\end{proof}
\begin{remark}
\textbf{构造动机与条件作用（重构新增）。}每层不可数，使可数次例外不能删尽该层；两个分支分别只允许删有限点、可数点，产生“不可数部分迫使前层几乎全部进入闭包”的不对称传播。对连续实函数，实线可由可数个短区间覆盖，因此这种传播将每层压缩到同一个数值。两端点的尾部邻域再把相同数值传到 $a,b$。另一方面，局部开闭邻域和尾部的闭包缩小保证空间仍正则。这两个机制分别负责函数约束与拓扑分离，不能相互替代。
\end{remark}
% END SUPPLEMENT: id=v14-ch04-regular-twins-construction

% SUPPLEMENT: id=v14-ch04-twin-quotient-construction; parent=ch04-notes-separation-history; kind=omitted-proof+construction+motivation+hypotheses; origin=independent-derivation-with-attributed-construction
\begin{proposition}[把两点约束传播到整个空间]\label{proposition:ch04-twin-quotient-regular}
\textbf{补充结论（重构新增）。}设 $S$ 是 $T_1$ 正则空间，$a\ne b\in S$。取非空一族不交副本 $S_s$，记其中相应两点为 $a_s,b_s$，令 $Z$ 为这些副本的拓扑和，$A=\{a_s\}$、$B=\{b_s\}$、$C=Z\setminus(A\cup B)$。若有双射 $g:A\to C$，把 $B$ 合并为一类，把每对 $\{a_s,g(a_s)\}$ 合并为一类，赋予所得集合 $Y$ 商拓扑。则 $Y$ 仍 $T_1$ 正则。若 $S$ 中所有连续实函数均在 $a,b$ 同值，则 $Y$ 上所有连续实函数都为常数。
\end{proposition}
\begin{proof}
\textbf{第一步：商与单点闭性。}双射保证列出的类两两不交且覆盖 $Z$。记商映射 $q:Z\to Y$。每个副本开闭，$S$ 正则且 $T_1$，所以 $Z$ 也有同样性质。$A,B$ 都是闭离散子集，并且它们的\emph{任意子集}在 $Z$ 中闭：与每个副本的交至多为一个闭单点。每个等价类不是闭集 $B$ 就是两个点组成的闭集，所以每个 $Y$ 单点的原像闭。按商拓扑定义，$Y$ 为 $T_1$。

\textbf{第二步：构造饱和的开邻域与闭包控制。}任给 $y\in Y$ 和包含它的开集 $O$。置 $U=q^{-1}(O)$、$E=q^{-1}(\{y\})$，则 $U$ 是包含整个类 $E$ 的饱和开集。我们构造饱和开集 $W$ 及饱和闭集 $K$，使
\begin{align*}
E\subseteq W\subseteq K\subseteq U.
\end{align*}
这样 $q(W)$ 开、$q(K)$ 闭，且 $y\in q(W)\subseteq\overline{q(W)}\subseteq q(K)\subseteq O$，即可证明正则性。

先选开集 $W_0$，满足 $E\subseteq W_0$、$\overline{W_0}\subseteq U$，并满足
\begin{align*}
\overline{W_0}\cap A=W_0\cap A=E\cap A,
\qquad
\overline{W_0}\cap B=W_0\cap B=E\cap B.
\end{align*}
这可在各副本中分别实现。若 $E=B$，在每个副本中给 $b_s$ 选闭包包含于 $U\cap S_s$ 且避开 $a_s$ 的开邻域，再取并。若 $E=\{a_{s_0},g(a_{s_0})\}$，只在包含这两个点的至多两个副本内选开邻域，其闭包留在 $U$ 内并避开所有不属于 $E$ 的标记点；每个副本仅有 $a_s,b_s$ 两个标记点，正则性及单点闭性保证能做到。其他副本取空集。副本均开闭，所以这些逐副本的闭包控制也在 $Z$ 中成立。

对每个 $a_s\in U$，\emph{预先固定一个}开邻域 $N_s\subseteq S_s$，使
\begin{align*}
a_s\in N_s,\qquad\overline{N_s}\subseteq (U\cap S_s)\setminus\{b_s\}.
\end{align*}
递归定义
\begin{align*}
W_{n+1}=W_n\cup\bigcup\{N_s:g(a_s)\in W_n\},
\qquad W=\bigcup_{n=0}^{\infty}W_n.
\end{align*}
每一步所需的 $N_s$ 都存在，因为 $W_n\subseteq U$ 且 $U$ 饱和，$g(a_s)\in W_n$ 蕴含 $a_s\in U$。各 $W_n$ 以及 $W$ 开，并包含于 $U$。

$W$ 饱和。若 $g(a_s)\in W$，它属于某个 $W_n$，则 $a_s\in N_s\subseteq W_{n+1}$。反过来，若 $a_s\in W$，要么它在 $W_0$ 中，此时由初始条件它属于 $E$，其伙伴也在 $E\subseteq W_0$；要么它因加入 $N_s$ 才进入，此时加入规则已经要求 $g(a_s)$ 在前一步中。其他副本的邻域不会加入 $a_s$。类 $B$ 也保持整体包含或整体排除，因为初始时如此，所有新增 $N_s$ 的闭包都避开 $b_s$。

关键是 $W$ 的闭包仍留在 $U$ 中。每个副本 $S_s$ 内，最终只可能出现 $W_0\cap S_s$ 与预先固定的\emph{一个} $N_s$ 的并，或只有 $W_0\cap S_s$。因此该副本内的闭包至多是两个已受控闭包的并，不会发生任意无限并的闭包逃出 $U$ 的问题。于是
\begin{align*}
L=\overline W\subseteq U,
\qquad L\cap A=W\cap A,
\qquad L\cap B=W\cap B\in\{\varnothing,B\}.
\end{align*}
其中 $A$ 的等式来自初始闭包避开未选中的 $a_s$，以及只有包含 $a_s$ 的 $N_s$ 才可能额外加入该点。

现在令
\begin{align*}
K=L\cup\{a_s:g(a_s)\in L\}.
\end{align*}
后加的集合是闭离散集 $A$ 的子集，所以闭，$K$ 为闭集。它包含于 $U$，因为 $L\subseteq U$ 且 $U$ 饱和。它也饱和：若 $g(a_s)\in L$，其伙伴按定义已加入；若 $a_s\in L$，则 $a_s\in W$，由 $W$ 饱和得 $g(a_s)\in W\subseteq L$。新增 $a_s$ 的伙伴本来也在 $L$ 中，且 $B$ 已整体包含或排除。故 $K$ 满足所需条件，正则性证明完成。

\textbf{第三步：连续函数恒定。}设 $h:Y\to\mathbb R$ 连续，记 $y_B$ 为整个 $B$ 所对应的商空间点，即任意 $q(b_s)$ 的共同值。$h\circ q$ 在每个副本 $S_s$ 上连续，所以由两点同值假设，
\begin{align*}
h(q(a_s))=h(q(b_s))=h(y_B).
\end{align*}
每个 $c\in C$ 都等于某个 $g(a_s)$，又与 $a_s$ 属于同一类，故它在 $h\circ q$ 下也取共同值 $h(y_B)$。$A,B,C$ 覆盖 $Z$，所以 $h$ 在全部 $Y$ 上恒定。
\end{proof}
\begin{remark}
\textbf{动机与逐条件说明（重构新增）。}把所有 $b_s$ 合并，使每个 $a_s$ 的函数值都被迫等于这个共同点的函数值；再把每个剩余点配给一个 $a_s$，约束就覆盖全部空间。双射使等价类确实构成划分。正则性的核心不是笼统地说“正则空间的商仍正则”——该说法一般不成立——而是这里每个副本只有一个可新加入的 $a_s$ 邻域，使闭包只需处理有限并；未饱和的闭包随后只需补入闭离散集的部分点。这些特殊结构不可省略。
\end{remark}
% END SUPPLEMENT: id=v14-ch04-twin-quotient-construction

% SUPPLEMENT: id=v14-ch04-regular-only-constant-proof; parent=ch04-continuous-functions-normality-intro; kind=omitted-proof+construction+motivation+hypotheses; origin=independent-derivation
\begin{corollary}\label{corollary:ch04-regular-only-constant}
\textbf{原存在性断言的补证结论（重构新增）。}存在至少含两个点的正则空间 $Y$，使每个连续实值函数及每个连续复值函数都为常数。
\end{corollary}
\begin{proof}
取\cref{proposition:ch04-regular-twins}中的 $S,a,b$，令 $T=S\setminus\{a,b\}$，固定一个 $t_0\in T$。以下用有限词直接构造上一命题所需的指标族与双射，无须把基数等式当作未验证的选择步骤。

令 $J=T^{<\omega}$ 为 $T$ 上全部有限词的集合，包括空词 $\epsilon$。以 $J$ 为指标取 $S$ 的不交副本；则 $A=\{a_s:s\in J\}$，剩余部分 $C$ 可识别为 $J\times T$。每个非空词有唯一的“前缀加末字母”分解，所以 $J\setminus\{\epsilon\}$ 与 $J\times T$ 之间有典范双射。再定义 $\theta:J\to J\setminus\{\epsilon\}$：
\begin{align*}
\theta(t_0^n)=t_0^{n+1}\quad(n\ge0),\qquad
\theta(s)=s\quad\text{若 }s\text{ 不全由 }t_0\text{ 组成}.
\end{align*}
这里 $t_0^0=\epsilon$。在全由 $t_0$ 组成的词链上它向后移一位，在其他词上不动，故确为双射。将 $a_s$ 送到 $\theta(s)$ 的前缀与末字母所对应的 $C$ 中点，便得到双射 $g:A\to C$。

\cref{proposition:ch04-twin-quotient-regular}给出正则商空间 $Y$，其所有连续实值函数恒定。$B$ 类与任意一个二点类不同，故 $Y$ 至少有两个点。若 $f:Y\to\mathbb C$ 连续，则实部、虚部都是连续实值函数，分别恒定，因而 $f$ 恒定。这补齐了\cref{section:ch04-continuous-maps}及本节注记中的存在性断言。
\end{proof}
\begin{remark}
\textbf{结论的意义（重构新增）。}此空间满足本书的 $T_1$ 正则性，所以任意两点仍有不交开邻域，但连续实函数却无法辨别它们。它不可能完全正则：若取不同点 $x,y$，闭单点 $\{y\}$ 与 $x$ 本应能由取值为 $0,1$ 的连续函数分离，与恒定性矛盾。这里“所有函数恒定”仅指连续实或复值函数，并非任意集合映射，也不把较弱的平凡拓扑例子冒充为正则例子。
\end{remark}
% END SUPPLEMENT: id=v14-ch04-regular-only-constant-proof
\end{document}
